%!TEX root = ../thesis.tex
% ******************************* Thesis Appendix A ****************************
\chapter{Simulation Hyperparameters}
\label{app:hyperparameters}

\noindent This appendix collects the numerical settings of every computational
experiment in the results chapters (Chapters~\ref{ch:model}--\ref{ch:conclusion}), so
that each reported figure and table can be regenerated exactly. The values are those of
the reference implementation, and together with the pseudocode of
Appendix~\ref{app:pseudocode} and the reference-code snapshots of
Appendices~\ref{app:code} (the \texttt{mnn-torch} package, for the spiking classifiers of
Chapter~\ref{ch:bio}) and~\ref{app:code-nn} (the \texttt{siox-eligibility-trace} package,
for the reinforcement-learning results of Chapters~\ref{ch:learning}--\ref{ch:conclusion})
they are sufficient to \emph{regenerate} the computational results without the source
repository by re-running the reference implementation; the per-cell \emph{outputs} of every swept
figure are additionally tabulated in Appendix~\ref{app:sourcedata}, so that each figure
can also be redrawn directly. They
are grouped by the quantity they parameterise: the fitted device model of
Chapter~\ref{ch:model} (Table~\ref{table:A0}); the device transient and the eligibility
trace it realises (Table~\ref{table:A1}); the shared learning-rule and neuron constants
(Table~\ref{table:A2}); the per-experiment task settings of the reinforcement-learning
ladder (Table~\ref{table:A3}); the additional constants of the deep all-local network
(Table~\ref{table:A4}) and of the hybrid perception front-end (Table~\ref{table:A5});
the hardware-aware spiking classifier of Chapter~\ref{ch:bio} with its Poole--Frenkel
device model (Table~\ref{table:A6}); and the exploratory extensions of
Chapter~\ref{ch:conclusion} (Table~\ref{table:A7}). Symbols follow the definitions in
the Nomenclature and in the chapters cited. The seconds-scale time constants are the
fitted device laws of Chapter~\ref{ch:model} (Table~\ref{table:A0}); $\tau_{leak}$ is
the programmable retention knob, not a fitted constant, and its swept values are listed
per experiment.

\paragraph{Scope of computational reproducibility.} The reinforcement-learning results
of Chapters~\ref{ch:learning} and~\ref{ch:conclusion} are pure-\texttt{numpy}
simulations of the fitted device model and reproduce in full from this appendix. The
hardware-aware classifier of Chapter~\ref{ch:bio} (Table~\ref{table:A6}) additionally
requires the \abbrev{MNIST}/N-\abbrev{MNIST} image datasets and a surrogate-gradient
spiking library; the two biological-reward capstones
(Figures~\ref{fig:7w2}--\ref{fig:7w}) require the external human-\abbrev{EEG}
(ds003474) and dopamine (\abbrev{DANDI}~000351) recordings; and the device-model
figures of Chapter~\ref{ch:model} are fits to measured instrument data. These
data-dependent results are reproducible in method but not from the thesis alone, and
are flagged as such where they occur.

% ---------------------------------------------------------------------------
\begin{table}[htbp!]
    \caption[Fitted device-model coefficients (Chapter~\ref{ch:model}).]{Coefficients
    of the compact dual-phase transient model of Chapter~\ref{ch:model}, fitted to
    current transients measured on a single gold-contacted SiOx device over
    $0.7$--$1.1$\,V and generalised by the field-acceleration laws of
    Equations~\ref{eq:5.8}--\ref{eq:5.9} (Table~\ref{table:5c}). $V=|\text{bias}|$ in
    volts. These are the source of the seconds-scale constants used throughout
    Table~\ref{table:A1}; the fit is a diagnostic of the model form, and the
    device-to-device spread is characterised only for the ITO retention study
    ($n=53$).}
    \begin{center}
    \small
    \begin{tabularx}{\textwidth}{@{}l X l@{}}
    \toprule
    Symbol & Meaning & Value \\
    \midrule
    $\beta_{fill}$   & Compressed-exponential rise exponent (fixed)     & $2$ \\
    $k$              & Sequential trap-cascade stages ($\beta_{fill}\!\to\!k$) & $3$ \\
    $\tau_{r0},\,c_r$ & Rise-law coefficients, $\tau_r=\tau_{r0}e^{-c_rV}$ & $1.45\times10^{2}$\,s,\ \ $2.9$ \\
    $\tau_{d0},\,c_d$ & Decay-law coefficients, $\tau_d=\tau_{d0}e^{-c_dV}$ & $1.17\times10^{4}$\,s,\ \ $3.9$ \\
    $V_{max}$        & Saturated trap occupancy (normalised)            & $1.0$ \\
    fit error        & Mean transient MSE (per-voltage refit)           & $1.9\%$ \\
    $n$              & Devices in the ITO retention study               & $53$ \\
    \bottomrule
    \end{tabularx}
    \label{table:A0}
    \end{center}
    \vspace*{-\baselineskip}
\end{table}

% ---------------------------------------------------------------------------
\begin{table}[htbp!]
    \caption[Device-transient and eligibility-trace constants.]{Fitted device
    transient and eligibility-trace constants, shared across all reinforcement-learning
    experiments of Chapter~\ref{ch:learning}. The rise and decay laws are the Chapter~5 fit
    (Table~\ref{table:5c}), with $V=|\text{bias}|$ in volts; the per-stage rate is
    $\alpha = k/\tau_r(V)$.}
    \begin{center}
    \small
    \begin{tabularx}{\textwidth}{@{}l X l@{}}
    \toprule
    Symbol & Meaning & Value \\
    \midrule
    $\beta_{fill}$   & Compressed-exponential rise exponent (fixed)          & $2$ \\
    $k$              & Number of sequential trap-cascade stages              & $3$ \\
    $\tau_r(V)$      & Cascade rise time constant (s)                        & $1.45\times10^{2}\,e^{-2.9V}$ \\
    $\tau_d(V)$      & Space-charge decay time constant (s)                  & $1.17\times10^{4}\,e^{-3.9V}$ \\
    $\alpha$         & Per-stage charging rate (s\textsuperscript{-1})       & $k/\tau_r(V)$ \\
    $V_{max}$        & Trap-occupancy bound (normalised state)               & $1.0$ \\
    $V$              & Operating bias magnitude (V)                          & $0.9$ (single layer), $1.5$ (deep, T-maze) \\
    $\tau_{leak}$    & Retention / credit-assignment window (s)              & programmable; see Table~\ref{table:A3} \\
    $\beta_{leak}$   & Discharge dispersion exponent                         & $1$ (single-rate); $0.85$ / $0.54$ (measured) \\
    \bottomrule
    \end{tabularx}
    \label{table:A1}
    \end{center}
    \vspace*{-\baselineskip}
\end{table}

% ---------------------------------------------------------------------------
\begin{table}[htbp!]
    \caption[Shared learning-rule and neuron constants.]{Constants shared by every
    reinforcement-learning experiment: the three-factor update (Equation~\ref{eq:7.8}),
    the signed leak-dominant coincidence (Equation~\ref{eq:7.9}), and the leaky
    integrate-and-fire action neuron (Equation~\ref{eq:7.10}).}
    \begin{center}
    \small
    \begin{tabularx}{\textwidth}{@{}l X l@{}}
    \toprule
    Symbol & Meaning & Value \\
    \midrule
    $\lambda_d$      & Acausal (LTD-wing) weight of the signed coincidence   & $0.6$ \\
    $\tau_m$         & LIF membrane time constant (s)                        & $20\times10^{-3}$ \\
    $v_{th}$         & LIF firing threshold                                  & $1.0$ \\
    $v_{reset}$      & LIF reset potential                                   & $0.0$ \\
    $\sigma$         & Membrane-noise (exploration) level                    & $0.15$ (annealed where noted) \\
    ---              & Reward-baseline tracking rate                         & $0.02$ per trial \\
    $b_0$            & Reward-baseline initial value ($A$ actions)           & $1/A$ (chance) \\
    $w_{init}$       & Initial synaptic weight                               & $0.5$ \\
    $w_{max}$        & Weight (conductance) bound                            & $1.5$ (bandit / deep); $1.0$ (distal reward) \\
    \bottomrule
    \end{tabularx}
    \label{table:A2}
    \end{center}
    \vspace*{-\baselineskip}
\end{table}

% ---------------------------------------------------------------------------
\begin{table}[htbp!]
    \caption[Per-experiment reinforcement-learning task settings.]{Per-experiment
    settings for the Chapter~\ref{ch:learning} ladder. All published results use $20$ random
    seeds run as a batched ensemble, with percentile-bootstrap $95\%$ confidence
    intervals ($10^4$ resamples). Sweeps list the swept values; $D$ is the
    action--reward delay and $\tau_{leak}$ the retention. Integration step $\Delta t$
    is in seconds.}
    \begin{center}
    \footnotesize
    \begin{tabularx}{\textwidth}{@{}X l l l l l@{}}
    \toprule
    Experiment (figure / table) & Task size & $\Delta t$ & $\tau_{leak}$ (s) & $D$ (s) & $\eta$ \\
    \midrule
    Trace-level distal reward (Fig.~\ref{fig:7c}) & $N=8$ synapses & $0.1$ & $\{1,5,20\}$ & $\{1\text{--}100\}$ & $0.05$ \\
    Spiking distal reward (Fig.~\ref{fig:7d}) & $N=8$ synapses & $10^{-3}$ & $\{0.5,2,10\}$ & $\{1\text{--}40\}$ & $0.2$ \\
    Contextual bandit (Fig.~\ref{fig:7e}) & $2\times2$ & $5\times10^{-3}$ & $10$ & $5$ & $0.2$ \\
    Delay $\times$ retention (Fig.~\ref{fig:7l}) & $2\times2$ & $5\times10^{-3}$ & $\{0.5,2,10\}$ & $\{1,2,5,10,20\}$ & $0.2$ \\
    Scaling (Fig.~\ref{fig:7g}) & $2\times2$ to $12\times8$ & $5\times10^{-3}$ & $10$ & $2$ & $0.2$ \\
    Remedies (Fig.~\ref{fig:7g}) & $8\times8$ & $5\times10^{-3}$ & $10$ & $2$ & $0.2$ \\
    Sequential T-maze (Table~\ref{table:7a}) & $L=3$, $2$ arms & $5\times10^{-3}$ & $\{2,5,20\}$ & $4$ & $0.2$ \\
    Deep XOR, all-local (Fig.~\ref{fig:7m}) & $4\to H\to2$ & $5\times10^{-3}$ & $10$ & $5$ & $0.2$ \\
    Capstone, EEG / dopamine (Fig.~\ref{fig:7w}) & $4\to16\to2$ & $5\times10^{-3}$ & $10$ & $5$ & $0.2$ \\
    Array-scale faults (Fig.~\ref{fig:7hs}) & $4\to H\to2$ & $5\times10^{-3}$ & $10$ & $5$ & $0.2$ \\
    Hybrid perception (Table~\ref{table:7c}) & $4$-way, learned feat. & $5\times10^{-3}$ & $10$ & $2$ & $0.2$ \\
    \bottomrule
    \end{tabularx}
    \label{table:A3}
    \end{center}
    \vspace{0.3em}
    \noindent\footnotesize\emph{Additional per-experiment settings.} All spiking tasks
    use Poisson input at $200$\,Hz on the active line ($40$\,Hz for distractors) with a
    cue epoch of $1.0$\,s (bandit/deep) or $2.0$\,s (distal reward). Trial/episode
    budgets: trace-level $1500$; spiking $400$; bandit $600$ ($6000$ for the remedy
    runs); T-maze $800$ episodes with a $0.4$\,s dwell per state; deep XOR $2000$;
    capstone $3000$. The scaling convergence criterion is $\tfrac{1}{2}(1+1/A)$. The
    remedy sweep anneals the exploration noise $\sigma:0.45\to0.05$. The e-prop baseline
    on the T-maze uses $\eta=2000$ (its computed eligibility is $\sim\!10^4\times$
    smaller). The array-scale fault sweep runs $H\in\{8,32,128,512\}$ at stuck fractions
    $\{0,0.05,0.20,0.50\}$ with a temporal distractor at $3.0$\,s ($12$ seeds). The
    EEG-decoded reward (dataset ds003474) is a logistic decoder on the frontocentral
    reward-positivity ($150$--$500$\,ms window), cross-validated $5$-fold ($C=0.1$); the
    dopamine reward is decoded from DANDI~000351.
\end{table}

% ---------------------------------------------------------------------------
\begin{table}[htbp!]
    \caption[Deep all-local network constants.]{Additional constants of the deep
    two-layer all-local network (direct feedback alignment $+$ local homeostasis,
    Equation~\ref{eq:7.14}). The first block is common to the deep XOR, capstone, and
    array-scale-fault experiments; the second lists the two settings that differ in the
    capstone. Weights are signed ($\pm w_{max}$), initialised as independent Gaussians
    of the given per-layer scale.}
    \begin{center}
    \small
    \begin{tabularx}{\textwidth}{@{}l X l@{}}
    \toprule
    Symbol & Meaning & Value \\
    \midrule
    $F,\,A$              & Input lines, action neurons (XOR)                 & $4,\ 2$ \\
    $H$                  & Hidden neurons                                    & $8$ (deep XOR); $16$ (capstone) \\
    $\eta_h$             & Hidden-layer learning rate                        & $\eta$ (deep XOR); $3.0$ (capstone) \\
    $B_{fb}$             & Feedback-matrix scale (fixed random)              & $1.0$ \\
    $c_{bias}$           & Tonic bias drive to action neurons during cue     & $0.35$ \\
    $\rho^{\ast}$        & Homeostatic target firing fraction                & $0.35$ \\
    $\tau_{h}$           & Homeostatic activity-estimate time constant       & $200$ trials \\
    $\eta_{h}^{homeo}$   & Homeostatic gain (0 disables)                     & $0.1$--$1.0$ \\
    $s_1,\,s_2$          & Per-layer weight-init scales ($W_1,W_2$)          & $0.6,\ 0.35$ \\
    $\sigma_0\!\to\!\sigma_1$ & Annealed exploration noise                   & $0.15\to0.05$ \\
    \midrule
    \multicolumn{3}{@{}l}{\emph{Capstone overrides:}\quad $\eta_h=3.0$,\ \ $B_{fb}=2.0$,\ \ $c_{bias}=0.3$.} \\
    \bottomrule
    \end{tabularx}
    \label{table:A4}
    \end{center}
    \vspace*{-\baselineskip}
\end{table}

% ---------------------------------------------------------------------------
\begin{table}[htbp!]
    \caption[Hybrid perception front-end constants.]{Constants of the gradient-trained
    convolutional spiking perception front-end used by the hybrid agent
    (Table~\ref{table:7c}). The front-end is trained by backpropagation-through-time and
    frozen; its low-dimensional readout then drives the device-trace decision layer,
    whose settings are the hybrid row of Table~\ref{table:A3}.}
    \begin{center}
    \small
    \begin{tabularx}{\textwidth}{@{}l X l@{}}
    \toprule
    Symbol & Meaning & Value \\
    \midrule
    ---            & Input grating size (pixels)                    & $16\times16$ \\
    $C$            & Orientation classes ($0,45,90,135^{\circ}$)    & $4$ \\
    ---            & Additive input-noise standard deviation        & $0.9$ \\
    ---            & Convolutional stack                            & $\text{conv}_8$--pool--$\text{conv}_{16}$--pool--FC \\
    ---            & Unrolled time steps                            & $20$ \\
    $\beta$        & LIF membrane decay                             & $0.9$ \\
    ---            & Surrogate gradient                             & fast sigmoid \\
    ---            & Optimiser / learning rate                      & Adam, $2\times10^{-3}$ \\
    ---            & Training steps / batch size                    & $600$ / $128$ \\
    \bottomrule
    \end{tabularx}
    \label{table:A5}
    \end{center}
    \vspace*{-\baselineskip}
\end{table}

% ---------------------------------------------------------------------------
\begin{table}[htbp!]
    \caption[Hardware-aware spiking-classifier settings.]{Training and device settings
    for the hardware-aware spiking networks of Chapter~\ref{ch:bio} (Table~\ref{table:6a}).
    The device block is the Poole--Frenkel model (Equations~\ref{eq:6.34}--\ref{eq:6.39});
    conductance bounds and the regression slopes are the fit to the measured multistate
    SiOx data, with the intercepts and residual covariance obtained from the same fit.}
    \begin{center}
    \small
    \begin{tabularx}{\textwidth}{@{}l X l@{}}
    \toprule
    Symbol & Meaning & Value \\
    \midrule
    ---            & Fully-connected network (MSNN)                 & $784\to100\to10$ \\
    ---            & Convolutional network (MCSNN)                  & $\text{conv}_{12}$--pool--$\text{conv}_{64}$--pool--FC$_{10}$ \\
    ---            & Unrolled time steps                            & $10$ \\
    $\beta$        & LIF membrane decay                             & $0.95$ \\
    ---            & Surrogate gradient                             & fast sigmoid \\
    ---            & Optimiser / learning rate                      & Adam, $5\times10^{-4}$ \\
    ---            & Batch size                                     & $64$ \\
    ---            & Seeds (reported mean $\pm$ s.d.)               & $3$ \\
    \midrule
    $k_V$          & Spike-to-voltage read gain                     & $0.5$ \\
    $V_{ref}$      & Poole--Frenkel reference voltage (V)           & $0.25$ \\
    $G_{off}$      & Minimum conductance (S)                        & $6.90\times10^{-4}$ \\
    $G_{on}$       & Maximum conductance (S)                        & $3.45\times10^{-3}$ \\
    ---            & Poole--Frenkel regression slopes               & $[0.237,\ 1.091]$ \\
    ---            & Temperature (embedded in the $k_BT$ term)      & $20\,^{\circ}\text{C}$ \\
    $r_{th}$       & Homeostatic target firing rate                 & $0.8$ \\
    $n_H$          & Homeostatic regulariser exponent               & $1$ \\
    ---            & Stuck-device fractions                         & $\{0,0.05,0.20,0.50\}$ \\
    $\sigma_g$     & Device-to-device (lognormal) spread            & $0.5$ \\
    \bottomrule
    \end{tabularx}
    \label{table:A6}
    \end{center}
    \vspace*{-\baselineskip}
\end{table}

% ---------------------------------------------------------------------------
\begin{table}[htbp!]
    \caption[Exploratory-extension settings (Chapter~\ref{ch:conclusion}).]{Settings for
    the exploratory extensions of Chapter~\ref{ch:conclusion}. Each is a variant of the
    reinforcement-learning ladder built on the same eligibility gate, three-factor rule
    and constants (Tables~\ref{table:A1}--\ref{table:A2}); only the departures are
    listed. All use $20$ random seeds with bootstrap $95\%$ confidence intervals
    ($12$ seeds for the array-scale case), $\Delta t = 5\times10^{-3}$\,s and
    $\eta = 0.2$ unless noted.}
    \begin{center}
    \footnotesize
    \begin{tabularx}{\textwidth}{@{}X l l X@{}}
    \toprule
    Experiment (figure) & $\tau_{leak}$ (s) & $D$ (s) & Departures from the ladder \\
    \midrule
    Long-horizon credit (Fig.~\ref{fig:8a}) & $10$ & $2$ & multi-step trajectory, length $L$ swept; $12$ seeds \\
    Reversal learning (Fig.~\ref{fig:8b}) & $10$ & $2$ & $2$ phases, reward contingency swapped mid-run; $3000$ trials \\
    Multi-timescale (Fig.~\ref{fig:8c}) & measured spread & $\{2,8\}$ & per-synapse $\tau$ drawn from the $n=53$ ITO fits; $S{=}4,A{=}2$ \\
    Vector timer (Fig.~\ref{fig:8d}) & $10$ & --- & per-stage cascade readout ($k{=}3$); two intervals $t_A{=}3$, $t_B{=}15$\,s; $\eta=0.3$, $4000$ trials \\
    WM $+$ STC, one device (Fig.~\ref{fig:8g}) & $\{20\text{--}40\}$ & $3$ & one substrate serves working memory and eligibility; $2500$ trials, $V{=}1.5$ \\
    Device-native TD (Fig.~\ref{fig:8h}) & $10$ & $2$ & actor--critic; discount $\gamma=e^{-\Delta t_{\text{state}}/\tau_{leak}}$; $L{=}5$, $2000$ episodes \\
    \bottomrule
    \end{tabularx}
    \label{table:A7}
    \end{center}
    \vspace*{-\baselineskip}
\end{table}

% ---------------------------------------------------------------------------
\begin{table}[htbp!]
    \caption[Reported final results with confidence intervals.]{Headline final
    performance of the reinforcement-learning results, as reported in
    Chapter~\ref{ch:learning}: mean over $20$ seeds with the percentile-bootstrap $95\%$
    confidence interval, and (where applicable) the number of seeds reaching criterion.
    These are the values plotted in the cited figures; they are collected here so a
    reader can check a reproduction against the exact numbers. Chance is $0.5$ for the
    two-choice tasks.}
    \begin{center}
    \footnotesize
    \begin{tabularx}{\textwidth}{@{}X l l l@{}}
    \toprule
    Result (figure) & Condition & Final reward rate [95\% CI] & Seeds solved \\
    \midrule
    Shallow DMS (Fig.~\ref{fig:7sdms}) & device trace     & $1.00\ [1.00,1.00]$ & --- \\
                                       & abstract kernel  & $0.86\ [0.85,0.87]$ & --- \\
                                       & no-trace control & $0.50\ [0.48,0.51]$ & --- \\
    \midrule
    Deep DMS, distractor (Fig.~\ref{fig:7dms}) & DFA $+$ homeostasis & $1.00\ [1.00,1.00]$ & $20/20$ \\
                                       & DFA, no homeostasis & $0.58\ [0.47,0.69]$ & $4/20$ \\
                                       & no-trace control    & $0.54\ [0.47,0.61]$ & $4/20$ \\
    \midrule
    Probabilistic selection (Fig.~\ref{fig:7q}) & shallow (device) & $0.86$; choose-A $0.91$, avoid-B $1.00$ & --- \\
                                       & no-trace control & $0.51$ & --- \\
    \bottomrule
    \end{tabularx}
    \label{table:A8}
    \end{center}
    \vspace*{-\baselineskip}
\end{table}

\FloatBarrier
